% GENERATED FILE. Edit canonical lessons, then run npm run generate:books.
% canonical-source-hash: fnv1a64:eb9f453e1ae4a529
% canonical-lessons: SA-C268-service-recall, SA-C268-vittakosah, SA-C268-post-health-location, SA-C268-bank-police-location, SA-C268-service-location-practice, SA-C268-service-location-retrieval

\chapter{Four services on one small map}
\label{ch:sa-service-locations}
\hlchaptermodality{268}

\begin{chapteropening}
\textbf{By the end of this chapter:} \emph{I can ask where four familiar public services are and answer here or there.}

Complete chapter 268 by independently asking and answering where four familiar service referents are.
\end{chapteropening}

\section[(service-word recall)]{Three service words are already yours}

\label{lesson:SA-C268-service-recall}



\begin{quote}
Say “here,” “there,” “where?”, and “is.” (\emph{Atra; tatra; kutra; asti.})
\end{quote}

\subsection*{Your turn}
\begin{itemize}
\raggedright
  \item \emph{Recall:} \textbf{\sk{पत्रालयः}} — \emph{patrālayaḥ} — a post office
  \item \emph{Recall:} \textbf{\sk{चिकित्सालयः}} — \emph{cikitsālayaḥ} — a hospital
  \item \emph{Recall:} \textbf{\sk{आरक्षकः}} — \emph{ārakṣakaḥ} — a policeman
\end{itemize}

\begin{grammarlens}[title={What you've built}]
The post, health, and police referents were already in the book. You also already know how to ask where something is. One missing service word comes next.
\end{grammarlens}

\begin{tcolorbox}[breakable,colback=teal!4,colframe=teal!35!black,title={Before you move on}]
Say the three service words again without looking.
\end{tcolorbox}
\section[\texorpdfstring{\sk{वित्तकोषः} (vittakoṣaḥ)}{वित्तकोषः (vittakoṣaḥ)}]{\sk{वित्तकोषः} (vittakoṣaḥ) — a financial bank}

\label{lesson:SA-C268-vittakosah}



\begin{quote}
Recall the Sanskrit for a riverbank, then the post office, hospital, and policeman. (\emph{Tīram; patrālayaḥ; cikitsālayaḥ; ārakṣakaḥ.})
\end{quote}

\subsection*{You'll want to know: \sk{वित्तकोषः}}
\textbf{\sk{वित्तकोषः}} — \emph{vittakoṣaḥ} — “a financial bank.” This is the institution that handles money, \textbf{not the riverbank} \textbf{\sk{तीरम्}} that you already know.

The modern financial sense is explicitly attested in the \href{https://kosha.sanskrit.today/glossary/places-bank}{Kosha Sanskrit Today “Places — Bank” glossary} and in the \href{https://sanskritdocuments.org/dict/dictallcheck.pdf}{Sanskrit Documents combined dictionary}. This lesson teaches only that sourced financial-institution sense.

\begin{cousinweb}
\textbf{\sk{वित्तकोषः}} is transparently formed from \textbf{\sk{वित्त}}, “wealth, money,” and \textbf{\sk{कोष}}, “treasury, store.” The compound names the financial institution; it does not change the older water-edge meaning of \textbf{\sk{तीरम्}}.
\end{cousinweb}

\subsection*{Your turn}
\begin{itemize}
\raggedright
  \item \emph{Say it:} \textbf{\sk{वित्तकोषः}} — a financial bank
  \item \emph{Contrast:} \textbf{\sk{वित्तकोषः}} — financial bank; \textbf{\sk{तीरम्}} — riverbank
\end{itemize}

\begin{tcolorbox}[breakable,colback=teal!4,colframe=teal!35!black,title={Before you move on}]
What is the sourced word in this lesson for a financial bank?
\end{tcolorbox}
\section[Practice]{One small four-service location exchange}

\label{lesson:SA-C268-post-health-location}



\begin{quote}
Say “financial bank,” then “here,” “there,” “where?”, and “is.”
\end{quote}

\subsection*{You'll want to know: four complete turns}
Keep the familiar question frame whole and change only its first word:

\begin{itemize}
\raggedright
  \item \textbf{\sk{पत्रालयः} \sk{कुत्र} \sk{अस्ति}? \sk{पत्रालयः} \sk{तत्र} \sk{अस्ति।}} — Where is the post office? The post office is there.
  \item \textbf{\sk{चिकित्सालयः} \sk{कुत्र} \sk{अस्ति}? \sk{चिकित्सालयः} \sk{अत्र} \sk{अस्ति।}} — Where is the hospital? The hospital is here.
  \item \textbf{\sk{वित्तकोषः} \sk{कुत्र} \sk{अस्ति}? \sk{वित्तकोषः} \sk{तत्र} \sk{अस्ति।}} — Where is the financial bank? The financial bank is there.
  \item \textbf{\sk{आरक्षकः} \sk{कुत्र} \sk{अस्ति}? \sk{आरक्षकः} \sk{अत्र} \sk{अस्ति।}} — Where is the policeman? The policeman is here.
\end{itemize}

These four complete turns are the whole pattern. Do not use them to infer a new declension or a broader directions system.

\subsection*{Your turn}
Ask where the post office is and answer “there.” Then ask where the hospital is and answer “here.”

\begin{tcolorbox}[breakable,colback=teal!4,colframe=teal!35!black,title={Before you move on}]
What stays fixed in all four questions? (\emph{Kutra asti?})
\end{tcolorbox}
\section[Practice]{The second pair, from memory}

\label{lesson:SA-C268-bank-police-location}



\begin{quote}
Say “financial bank” and “policeman.”
\end{quote}

\subsection*{Your turn}
\begin{itemize}
\raggedright
  \item \emph{Ask and answer:} Where is the financial bank? The financial bank is there.
  \item \emph{Ask and answer:} Where is the policeman? The policeman is here.
\end{itemize}

Check only after speaking: \textbf{\sk{वित्तकोषः} \sk{कुत्र} \sk{अस्ति}? \sk{वित्तकोषः} \sk{तत्र} \sk{अस्ति।} \sk{आरक्षकः} \sk{कुत्र} \sk{अस्ति}? \sk{आरक्षकः} \sk{अत्र} \sk{अस्ति।}}

\begin{grammarlens}[title={What you've built}]
You have now retrieved both halves of the four-service exchange. Nothing here adds directions, a police station, or a service transaction.
\end{grammarlens}

\begin{tcolorbox}[breakable,colback=teal!4,colframe=teal!35!black,title={Before you move on}]
Ask the two questions once more, without looking.
\end{tcolorbox}
\section[(independent service locations, first …]{Independent service locations: round one}

\label{lesson:SA-C268-service-location-practice}



\begin{quote}
Recall the four small location pieces: here, there, where, is.
\end{quote}

\subsection*{Your turn: one small map}
No Sanskrit model is shown before you answer.

Ask where each of these is: the post office, hospital, financial bank, and policeman. Answer “there” for the post office and bank, and “here” for the hospital and policeman.

\begin{grammarlens}[title={What you've built: boundary}]
You can complete one small four-service location round. You have not yet learned directions, transactions, emergency requests, or institutional systems.
\end{grammarlens}

\begin{tcolorbox}[breakable,colback=teal!4,colframe=teal!35!black,title={Before you move on}]
Name the four referents again without looking.
\end{tcolorbox}
\section[(independent service locations, second …]{Independent service locations: round two}

\label{lesson:SA-C268-service-location-retrieval}



\begin{quote}
Contrast the financial bank with the riverbank one last time.
\end{quote}

\subsection*{Your turn: a fresh arrangement}
No Sanskrit model is shown before you answer.

Ask where each of the same four referents is. This time answer “here” for the post office and financial bank, and “there” for the hospital and policeman.

\begin{grammarlens}[title={What you've built: boundary}]
You can retrieve four service referents and one bounded location exchange. This is not evidence of directions, postal or banking transactions, medical or emergency communication, police-station vocabulary, human validation, learner success, or exam readiness.
\end{grammarlens}

\begin{tcolorbox}[breakable,colback=teal!4,colframe=teal!35!black,title={Before you move on}]
Name the four services, then name the one exchange you can now complete with each.
\end{tcolorbox}
